% sections/model.tex -- system model, binding soundness, failure classes, adversaries, scope.
% Terms follow binding/CONTRACT.md. Fig. 1 (architecture) and Fig. 2 (views) are TikZ (W25).
\section{System Model and Binding Soundness}\label{sec:model}

\begin{figure}[t]
\centering
\begin{tikzpicture}[font=\scriptsize,
  box/.style={draw, rounded corners=2pt, align=center, inner sep=2.5pt, minimum height=0.9cm},
  adv/.style={draw, dashed, rounded corners=2pt, align=center, inner sep=2pt, fill=red!4},
  cls/.style={fill=red!12, draw=red!60!black, text=red!50!black, rounded corners=1pt,
     inner sep=1.2pt, font=\scriptsize\bfseries},
  >={Stealth[length=3.5pt]}]
\node[box] (C) at (0,0) {Client $\Cl$\\ key $k$, token $T$};
\node[box] (G) at (2.45,0) {Gateway\\ (+ PDP $\PDP$)};
\node[box, minimum width=2.5cm] (E) at (5.75,0) {Executor $\Exe$\\ override $\to$ router\\ $\to$ handler};
\node[box] (V) at (2.45,-1.75) {Verifier $\Ver$\\ Alg.~\ref{alg:verifier} on $\view(r)$};
\draw[->] (C) -- node[above]{$r,T,\pi$} (G);
\draw[->] (G) -- node[above]{$\tau_D(r)$} (E);
\draw[->] ([xshift=-5pt]G.south) -- node[left, align=right]{edge: view\\ headers} ([xshift=-5pt]V.north);
\draw[->] ([xshift=5pt]V.north) -- node[right, align=left]{allow,\\ \hdr{X-Verified-*}} ([xshift=5pt]G.south);
\draw[->, dashed] ([xshift=8pt]E.south) |- node[pos=0.22, right, align=left]{origin:\\ $\Exe$ calls $\Ver$} (V.east);
\node[above=0pt of E, font=\scriptsize] {$\exe(r)=(m_x,\mathit{route}_x,\op_x)$};
\node[cls, anchor=south] at ([xshift=12pt]E.north east) {G1};
\node[cls] at (4.05,-0.55) {G2};
\node[cls, anchor=south] at ([xshift=-4pt]G.north west) {G3};
\node[adv, anchor=north west] (AL) at (-0.75,-1.15) {$\Aleak$: replays one\\ leaked $(T,\pi)$};
\node[adv, anchor=north west] (AC) at (-0.75,-2.0) {$\Aclient$: signs\\ any request};
\node[align=left, anchor=north west, font=\scriptsize, text width=8.0cm] at (-0.75,-2.8) {%
\textbf{\textcolor{red!50!black}{G1}} query, body or override $\notin\Cov$ select $\op_x$;
\textbf{\textcolor{red!50!black}{G2}} $\view(r)$ and $\Exe$'s parse differ;
\textbf{\textcolor{red!50!black}{G3}} client sets a header a later hop trusts.};
\end{tikzpicture}
\Description{System and threat architecture: a client with key and token sends a request with a proof to a gateway that may host a policy decision point; the gateway forwards it to an executor made of override middleware, router and handler. Below the gateway, a verifier checks the proof against the view headers (edge placement); in origin placement the executor calls the verifier itself (dashed). Red tags mark where the three failure classes arise: G3 at the gateway boundary, G2 between the verifier's view and the executor's parse, G1 inside the executor. Two adversaries are shown: one who replays a leaked token and proof, and a legitimate client who signs any request.}
\caption{System and threat architecture. $\Ver$ checks $\pi$ against its view
$\view(r)$; $\Exe$ runs $\exe(r)$ on the request $\tau_D(r)$ the intermediaries
deliver. Edge placement reaches $\Ver$ through a gateway adapter; origin
placement calls it from $\Exe$ (dashed).}
\label{fig:arch}
\end{figure}

\subsection{System model}\label{sec:model:system}

A deployment $D$ consists of a client, a verifier, an optional policy decision
point, zero or more intermediaries and an executor (Figure~\ref{fig:arch}).
The client $\Cl$ holds a key $k$ and a token $T$ with $T.\jkt=\mathsf{th}(k)$,
the RFC~7638 thumbprint. For each request $r$ it signs a fresh proof
$\pi=(\htm,\htu,\ath,\jti,\iat)$. The \emph{canonical request} $r_\pi$ is the
request an honest client sends with $\pi$: method \htm, origin-form target the
path of \htu. Otherwise $r$ is an arbitrary byte string.

The verifier $\Ver$ maps the request it observes to a \emph{view}
$\view(r)=(m_v,u_v)$, taken from the gateway's own request line or from headers
such as \hdr{X-Forwarded-Uri} the gateway writes. $\Ver$ \emph{accepts} iff
Algorithm~\ref{alg:verifier} allows on that view; its comparison function
$\canon$ is a deployment parameter (RFC~9449 only recommends RFC~3986
normalisation). An optional PDP $\PDP$ maps the view and the subject to a
decision on a resource $\rho_\PDP(\view(r))$, e.g.\ a path-keyed policy.
The executor $\Exe$---router, override middleware and handler---receives
$\tau_D(r)$ and computes the \emph{executed operation}
$\exe(r)=(m_x,\mathit{route}_x,\op_x)$: the effective method after override,
the matched route with parameter assignment $\theta_x$, and the operation that
ran, including dispatch on the query or body. Two executions are the same
operation iff they agree on $\hat\exe(r)=(\op_x,\theta_x)$. Ground truth is
$\op(\pi)=\hat\exe(r_\pi)$ for $r_\pi$ sent directly to $\Exe$ without
intermediaries or verification.

\subsection{Binding soundness}\label{sec:model:soundness}

Let $\Acc_D(\pi)$ be the set of requests $r$ that the verifier accepts under
$\pi$ in $D$ and for which the executor runs an operation.

\begin{definition}[Binding soundness]\label{def:soundness}
$D$ is \emph{binding-sound} if, for every proof $\pi$ and all
$r,r'\in\Acc_D(\pi)$, $\hat\exe(r)=\hat\exe(r')$; it is \emph{binding-correct}
if moreover $\hat\exe(r)=\op(\pi)$ for every $r\in\Acc_D(\pi)$.
\end{definition}

\begin{definition}[Binding resolution]\label{def:resolution}
$\res_D(\pi)=\lvert\{\hat\exe(r) : r\in\Acc_D(\pi)\}\rvert$.
\end{definition}

$D$ is binding-sound iff $\res_D(\pi)\le 1$ for all $\pi$. In a functioning
deployment $r_\pi\in\Acc_D(\pi)$, so one accepted $r$ with
$\hat\exe(r)\neq\op(\pi)$ witnesses $\res_D(\pi)\ge 2$
(Figure~\ref{fig:views}). Our oracle searches for such witnesses; since it
observes finitely many requests, every reported resolution is a lower bound.
Binding soundness is weaker than message integrity: requests may differ in any
component provided the executed operation does not. It is the property
RFC~9449 implicitly relies on when it signs ``just enough''.

\begin{figure}[t]
\centering
\begin{tikzpicture}[font=\scriptsize, >={Stealth[length=3.5pt]},
  dot/.style={circle, fill, inner sep=1.1pt},
  set/.style={draw, rounded corners=6pt, inner sep=4pt}]
\node[dot, label={[font=\scriptsize]left:$r_\pi$}] (r0) at (0,0.55) {};
\node[dot, label={[font=\scriptsize]left:$r_1$}] (r1) at (0,0) {};
\node[dot, label={[font=\scriptsize]left:$r_2$}] (r2) at (0,-0.55) {};
\node[set, fit=(r0)(r1)(r2), minimum width=1.3cm, label={[font=\scriptsize]above:$\Acc_D(\pi)$}] (A) {};
\node[font=\scriptsize, align=center] at (0,-1.25) {$\Ver$: $\canon(\view(r))\equiv(\htm,\htu)$};
\node[dot, label={[font=\scriptsize]right:$\op(\pi)$, e.g.\ \texttt{deleteAccount}}] (o0) at (3.2,0.55) {};
\node[dot, label={[font=\scriptsize]right:\texttt{deleteStatement} (\GB)}] (o1) at (3.2,0) {};
\node[dot, label={[font=\scriptsize]right:\texttt{accountAction.close} (\GA)}] (o2) at (3.2,-0.55) {};
\node[font=\scriptsize, align=center] at (4.35,-1.25) {$\Exe$: $\hat\exe(r)=(\op_x,\theta_x)$};
\draw[->] (r0) -- (o0); \draw[->, red!60!black] (r1) -- (o1); \draw[->, red!60!black] (r2) -- (o2);
\node[font=\scriptsize] at (1.6,0.95) {$\hat\exe$};
\end{tikzpicture}
\Description{Two columns of dots. On the left, the set of requests the verifier accepts under one proof: the canonical request and two mutated requests, all of which map to the same verifier view. On the right, executed operations. The canonical request maps to the proof's operation; the two mutated requests map to two other operations, one through an interpretation differential and one through query dispatch. Three distinct operations means binding resolution three.}
\caption{Binding soundness. All requests in $\Acc_D(\pi)$ look the same to
$\Ver$ (verifier view $\view(r)$), but $\Exe$ executes $\hat\exe(r)$; here
$\res_D(\pi)=3$, witnessed by a \GB and a \GA request.}
\label{fig:views}
\end{figure}

\begin{definition}[Authorization soundness]\label{def:authz}
A deployment with $\PDP$ is \emph{authorization-sound} if, whenever $\Ver$
accepts $r$ and $\PDP$ permits action $a$ on $\rho=\rho_\PDP(\view(r))$, the
executor acts only on $\rho$ with action $a$.
\end{definition}

Binding soundness concerns which operation a \emph{proof} selects;
authorization soundness concerns which resource a \emph{decision} selects. A
binding-soundness finding always involves a proof the adversary did not
create; an authorization-soundness finding involves none, because the client
signs exactly what it sends. We report the two separately and never add them.

\subsection{Failure classes}\label{sec:model:classes}

Let $\Cov=\{\text{method},\text{target URI without query},T\}$ be the covered
components and $\mathsf{in}_\Exe(r)=(m^\Exe,p^\Exe)$ the method and path as
the executor parses them, before override. A witness falls into one or more
classes.
\textbf{\GA, under-binding by design:} $\Exe$ reads the covered components as
$\Ver$ does, $(m^\Exe,\canon(p^\Exe))=(m_v,\canon(u_v))$, yet the operation
differs, so something outside $\Cov$ selected it: a dispatching query
parameter, a body field (GraphQL \texttt{operationName}, JSON-RPC), or a
method override (\hdr{X-HTTP-Method-Override}, \texttt{\_method}).
\textbf{\GB, interpretation differential:} $\Ver$ and $\Exe$ disagree on a
covered component, $(m^\Exe,\canon(p^\Exe))\neq(m_v,\canon(u_v))$---dot
segments, encoded separators, fragments, matrix parameters, duplicate slashes,
case, prefix rewrites.
\textbf{\GC, binding propagation:} a field that a downstream component treats
as the result of an upstream check (\hdr{X-Original-URI},
\hdr{X-Verified-Sub}, \hdr{Client-Cert}) is chosen by the client. \GC can make
$\view(r)$ client-controlled, turning acceptance into a \GB witness, or change
the subject the executor believes was verified.

\subsection{Adversaries}\label{sec:model:adversaries}

\paragraph{\Aleak, proof-leak adversary.} \Aleak obtains one valid
$(T,\pi)$, e.g.\ from logs or an intermediary's request log, and can send
arbitrary bytes to the public origin. It does not learn $k$, cannot mint or
alter proofs and cannot break TLS. Its goal is, within the \iat window, an
executed $\hat\exe(r)\neq\op(\pi)$. This matches attacker A5 of the FAPI~2.0
attacker model, which ``can also read requests sent to the resource
server''~\cite{fapiam2025}, though A5 sees them only after processing. Replay
state limits \Aleak but not everywhere: RFC~9449 notes that single-use \jti
checks may be infeasible without shared state and frames them ``in the
context of the target URI''. The harness sends only the mutated request per
fresh proof, with the replay cache on; this models deployments whose \jti
check does not reject it (the canonical request went to another replica, or
the cache is keyed per URI).

\paragraph{\Aclient, malicious legitimate client.} A registered client with
its own $k$ and $T$ and consent for resources $R$ (\texttt{alice} owns
\texttt{acc-alice-1}). It can mint any proof. Its goal is that $\Ver$ and
$\PDP$ approve $\rho\in R$ while $\Exe$ acts on $\rho'\notin R$
(Definition~\ref{def:authz}).

\paragraph{Excluded and scope.} We exclude network attackers (TLS is
assumed), malicious operators of gateway, verifier or executor, malware
holding $k$, a compromised authorization server~\cite{hosseyni2024csf}, and
denial of service. Harness traffic is cleartext HTTP/1.1 to
\texttt{http://api.bench.test}: view extraction, URI comparison, routing,
override and header forwarding happen after TLS termination, neither
adversary needs to break TLS, and certificate binding behind a proxy reduces
to the header that our \GC primitives test. HTTP/2 front ends appear only via
\texttt{h2c}. Hypotheses \Hyp{1}--\Hyp{5} and their falsifiers are in
Appendix~\ref{app:hyp}.
